\section{Scoring the uncertainty: the error chain along \texorpdfstring{$\tau$}{tau}}\label{sec:fms}

The chain~\eqref{eq:chain0} scores the mean of a scheme. A DA scheme also reports
an uncertainty --- an ensemble spread, a Gaussian covariance, a set of samples ---
and the paper's subtitle asks whether that uncertainty is right. We score it with
a single proper score that reduces to the MSE at one end of its range, so that the
same chain applies.

\subsection{The operator of a scheme}
Work in per-channel standardised units, and let $\xv_1$ denote a target state
(a draw from the scheme's output, or the truth $\xv_1^\star$). On the linear path
\begin{equation}
  \xv_\tau = \tau\, \xv_1 + (1-\tau)\, \xv_0, \qquad \xv_0 \sim \mathcal{N}(0, I),\quad \tau \in [0, 1),
\end{equation}
any output distribution $q$ of a scheme defines the operator
\begin{equation}\label{eq:psi}
  \Psi_q(\xv_\tau, \tau) \;=\; \E_q[\xv_1 \mid \xv_\tau] \;=\; \Psi_{\mathrm{mean}} + \Psi_{\mathrm{anom}},
\end{equation}
the denoiser of $q$ at noise level $\tau$, split into the scheme's mean
$\Psi_{\mathrm{mean}} = \Psi_q(\cdot, 0)$ and an anomaly part that carries its
uncertainty. Table~\ref{tab:psi} gives $\Psi_q$ for the three kinds of output we
meet. A point estimate is a degenerate flow. A Gaussian has a closed form. An
ensemble, a learned sampler or a kernel density of members gives a
posterior-weighted average of its samples. Conditional flow matching learns
$\Psi$ directly~\todo{cite Lipman et al.\ 2023}; the general representation of
DA algorithms as flows is beyond this paper.

\begin{table}[t]
\centering\small
\caption{The operator $\Psi_q$ of each kind of scheme output ($\snr = \tau^2/(1-\tau)^2$; per element for the Gaussian form).}
\label{tab:psi}
\begin{tabular}{lll}
\toprule
schemes & output $q$ & $\Psi_q(\xv_\tau,\tau)$ \\
\midrule
Strong/Weak-4DVar, DirectUNet & point mass at $\hat m$ & $\hat m$ \quad ($\Psi_{\mathrm{anom}} \equiv 0$) \\
EnKF, ETKF, ETKS (Gaussian form) & $\mathcal{N}(m, v)$ & $m + \dfrac{\snr\, v}{1 + \snr\, v}\,(\xv_\tau/\tau - m)$ \\
CFM, SDA, ensemble (KDE) & samples $\{\xv_1^{(i)}\}$ & $\sum_i w_i(\xv_\tau)\, \xv_1^{(i)}$, \ $w_i \propto \mathcal{N}(\xv_\tau; \tau \xv_1^{(i)}, (1-\tau)^2 I)$ \\
hybrid DirectUNet $\to$ SDA3-fix & composition & $\Psi_{\mathrm{mean}}$ from DirectUNet, $\Psi_{\mathrm{anom}}$ from SDA \\
\bottomrule
\end{tabular}
\end{table}

\subsection{The score}
The flow-matching score of a scheme is the error of its operator against the
truth,
\begin{equation}\label{eq:fms}
  \FMS_\tau(q) \;=\; \E\big\|\Psi_q(\xv_\tau^\star, \tau) - \xv_1^\star\big\|^2,
  \qquad \xv_\tau^\star = \tau\, \xv_1^\star + (1-\tau)\, \xv_0 .
\end{equation}
It has six properties we use (proofs in Appendix~\ref{app:proofs}).
\begin{enumerate}[label=(\roman*),leftmargin=*,nosep]
  \item At $\tau = 0$ it is the mean-squared error of the scheme's mean.
  \item For $\tau \in (0,1)$ it is strictly proper: its expectation is minimised
    only by the true conditional distribution, because Gaussian smoothing is
    injective.
  \item For a Gaussian $\mathcal{N}(m, v)$ with error $e$, $\FMS_\tau = (e^2 +
    \snr\, v^2)/(1 + \snr\, v)^2$, minimised at $v = \E e^2$ for every $\tau$:
    the spread--skill identity.
  \item As $\tau \to 1$ it tends to a universal floor $1/\snr$ plus a term
    proportional to the Hyvärinen score.
  \item Integrated with weight $d\snr$ it gives the log score up to a constant
    (mismatched estimation) \todo{cite Verdú 2010}.
  \item A point estimate scores its MSE at every $\tau$.
\end{enumerate}
Small $\tau$ reads the mean; large $\tau$ reads the uncertainty, and in particular
penalises a confident ensemble whose truth falls outside it.

\subsection{The chain at every \texorpdfstring{$\tau$}{tau}}
Let $\Psi_p$ and $\Psi_{p_S}$ be the operators of the true conditional
distributions given $C$ and given $S$. Under the condition
of~\eqref{eq:chain0}, for every $\tau$,
\begin{equation}\label{eq:chaintau}
  \FMS_\tau(q) \;=\;
  \underbrace{\mathrm{mmse}_p(\tau)}_{\text{irreducible}}
  \;+\; \underbrace{\E\|\Psi_p - \Psi_{p_S}\|^2}_{R_\tau(S):\ \text{restriction}}
  \;+\; \underbrace{\E\|\Psi_q - \Psi_{p_S}\|^2}_{D_\tau(q \,\|\, p_S):\ \text{misspecification + approximation}} ,
\end{equation}
with $\mathrm{mmse}_p(\tau) = \E\|\Psi_p - \xv_1^\star\|^2$. The proof repeats
the tower-property argument of~\eqref{eq:chain0} conditioning on $(\xv_\tau, S)$
instead of $S$: $\Psi_{p_S} = \E[\Psi_p \mid \xv_\tau, S]$, and $\Psi_q - \Psi_{p_S}$
is a function of $(\xv_\tau, S)$ and of independent randomness. $D_\tau \ge 0$,
with equality if and only if $q = p_S$. At $\tau = 0$, \eqref{eq:chaintau}
reduces to~\eqref{eq:chain0}. Three consequences follow.
\begin{itemize}[leftmargin=*]
  \item \textbf{Differences are exact.} $\mathrm{mmse}_p(\tau)$ is unknown but
    common to all schemes, so differences of $\FMS_\tau$ between schemes are
    exact differences of operator divergences.
  \item \textbf{A filter/smoother identity along $\tau$.} If a filter and a smoother computed on the same ensemble were both exact,
    their score gap would equal $\E\|\Psi_F - \Psi_S\|^2$ at every $\tau$. A gap
    that does not match flags a misspecification. At $\tau = 0$ this is the
    identity $\mathrm{MSE}_F - \mathrm{MSE}_S = \mathrm{var}_F - \mathrm{var}_S =
    \E\|m_F - m_S\|^2$ that we test in Sect.~\ref{sec:q1}.
  \item \textbf{Restriction vs misspecification, by eye.} A restriction raises the
    whole $\FMS_\tau$ curve, and the scheme stays calibrated for its information. A
    misspecification makes the scheme confident about the wrong thing, and bends
    the curve upward at large $\tau$.
\end{itemize}

\subsection{Derived diagnostics}
\begin{itemize}[leftmargin=*]
  \item \textbf{Calibration loss}: the share of $\FMS_\tau$ removed by the best
    scalar rescaling $c\,v$ of the reported variance, with its optimum
    $c^\star(\tau)$. A scheme can be calibrated in aggregate and still lose at
    large $\tau$, if its spread is right on average but in the wrong place.
  \item \textbf{Shape gain}: the score of a kernel density of the members against a
    Gaussian of the same variance; it reads non-Gaussian structure.
  \item \textbf{Pooled spread/skill} $\sqrt{\E v / \E e^2}$ ($\approx 0.984$ for a
    calibrated 30-member ensemble). The ratio of per-window mean spread to
    per-window mean RMSE is biased low and is not used as a calibration measure.
\end{itemize}

\paragraph{Relation to CRPS and the log score.}
The chain decomposition is not unique to $\FMS$: the expected log score splits
the same way, through Kullback--Leibler divergences, and calibration--resolution
decompositions exist for every proper score \todo{cite Bröcker 2009; Gneiting \&
Raftery 2007}. We use $\FMS_\tau$ for three reasons:
\begin{itemize}[leftmargin=*,nosep]
  \item it is computable for samples, Gaussians and point estimates alike, where
    the log score is undefined for ensembles and point masses;
  \item it is resolved over $\tau$, which separates the mean from the uncertainty;
  \item it equals the MSE at $\tau = 0$, which keeps continuity with DA practice.
\end{itemize}
The CRPS remains a useful marginal check and is reported alongside.

\figplaceholder{Fig.~2 -- the chain along $\tau$}{Synthetic Gaussian example: $\FMS_\tau$ of the true posterior, of a restricted one (curve raised, same shape) and of a misspecified one (bent upward at large $\tau$); the three terms of~\eqref{eq:chaintau} shaded.}
